\documentclass[12pt]{article}

% Packages
\usepackage[utf8]{inputenc}
\usepackage{amsmath, amssymb}
\usepackage{geometry}
\usepackage{tikz} % For drawing
\usepackage{tikz-3dplot} % For 3D plotting
\usepackage{float} % For figure placement
\geometry{margin=1in}


\usepackage{enumerate}
\geometry{margin=1in}

\title{Calculus II Week 1 Discussion Response}
\author{Elijah}
\date{\today}

\begin{document}

\maketitle

\subsection{Prompt}
Create an integral for calculating the volume of a sphere (disk method). Use it to "prove" the formula for the volume of a sphere.

\subsection{Response}
Rewriting \((x-h)^2 + (y-k)^2 = r^2\), we have
\[
(y-k)^2 = r^2 - (x-h)^2 \implies y-k = \sqrt{r^2 - (x-h)^2} \implies y = \sqrt{r^2 - (x-h)^2} + k,
\]
which gives the positive area of the circle. For the sake of calculating the volume of a sphere, we should assume constants \((h,k)=(0,0)\). Then, 
\[y=\sqrt{r^2-x^2}\]

Then, the volume of one disk perpendicular to the \(x\)-axis for some \(x\) will have radius \(r = \sqrt{r^2-x^2}\) and height \(dx\). So, the volume of one disk is \(\pi r^2dx\). Then, the sum of all disks from \(-r\) to \(r\) (and therefore the volume of a sphere with radius \(r\)) is \[V=\pi\int_{-r}^r r^2dx=\pi\int_{-r}^r\left[\sqrt{r^2-x^2}\right]^2dx=\pi\int_{-r}^{r}r^2-x^2dx\] 
\[=\pi\left[r^2x-\frac{x^3}{3}\Big|_{-r}^{r}\right]=\pi\left[\left(r^2(r)-\frac{r^3}{3}\right)-\left(r^2(-r)-\frac{(-r)^3}{3}\right)\right]\]
\[=\pi\left[\left(\frac{3r^3-r^3}{3}\right)-\left(\frac{-3r^3-(-r^3)}{3}\right)\right]=\pi\left[\left(\frac{3r^3-r^3}{3}\right)-\left(\frac{-3r^3+r^3}{3}\right)\right]\]

\[=\pi\left[\left(\frac{2r^3}{3}\right)-\left(\frac{-2r^3}{3}\right)\right]=\boxed{\frac{4\pi r^3}{3}}\]

\end{document}